
\chapter*{Introducción}
\addcontentsline{toc}{chapter}{Introducción}

Sin duda, la energía eléctrica es vital para el progreso de cada región, tanto social, económica como industrialmente. En Colombia, el acceso a este servicio está estrictamente regulado en términos de provisión, con el objetivo de asegurar que todos los usuarios cuenten con un suministro de energía constante, seguro y de alta calidad para vivir y trabajar~\cite{ref1}. Los operadores de red tienen de manera continua la responsabilidad de garantizar que la entrega de energía sea confiable para hogares, instituciones y negocios, evitando fallas en la prestación del servicio en cualquier tipo de edificación.

No obstante, esta tarea no es sencilla. La operación de los sistemas de distribución de energía se ve afectada por múltiples factores técnicos, ambientales y operativos, así como por la duración y frecuencia de las interrupciones del servicio. En este contexto, la Comisión de Regulación de Energía y Gas (CREG) adoptó los indicadores SAIDI (Índice de Duración Promedio de Interrupción del Sistema) y SAIFI (Índice de Frecuencia Promedio de Interrupción del Sistema) como los criterios oficiales para evaluar la calidad del servicio en las redes de distribución, mediante la Resolución CREG 015 de 2018 \cite{ref1}. Estos indicadores permiten evaluar el desempeño operativo y facilitan la implementación de esquemas de compensación y control por parte de la Superintendencia de Servicios Públicos Domiciliarios (SSPD) \cite{ref4}.

Sin embargo, la gestión efectiva de estos indicadores ha resultado compleja, dada la heterogeneidad de los eventos que originan las interrupciones y la limitada conectividad de la información almacenada en sistemas dispares como SCADA, OMS, GIS, CIS y bases de datos externas.\footnote{SCADA (\textit{Supervisory Control and Data Acquisition}): sistema de supervisión y adquisición de datos de la operación de la red. OMS (\textit{Outage Management System}): sistema de gestión de interrupciones. GIS (\textit{Geographic Information System}): sistema de información geográfica de activos. CIS (\textit{Customer Information System}): sistema de información de clientes.} Debido a que los registros operativos, climáticos y de mantenimiento se encuentran dispersos y fragmentados en múltiples plataformas, construir modelos analíticos capaces de pronosticar la frecuencia y el tiempo de duración de interrupciones eléctricas en distintas áreas o circuitos representa un desafío significativo.

Esta brecha limita la capacidad de las empresas para tomar decisiones oportunas, priorizar inversiones, mejorar la planificación del mantenimiento y responder adecuadamente a los requerimientos regulatorios. En consecuencia, resulta fundamental desarrollar herramientas apoyadas en ciencia de datos que permitan consolidar información histórica, identificar patrones relevantes y pronosticar la frecuencia y el tiempo de duración de interrupciones eléctricas.

Este trabajo presenta técnicas para la recopilación de datos provenientes de plataformas operativas internas y fuentes externas, así como el análisis de modelos estadísticos y de aprendizaje automático orientados al \textbf{pronóstico de la frecuencia y el tiempo de duración de interrupciones eléctricas} en una empresa de distribución de energía que opera en el departamento de Norte de Santander, en el marco de los indicadores regulatorios SAIDI y SAIFI. Al adoptar esta estrategia, se busca anticipar cuándo y con qué intensidad operativa pueden ocurrir las interrupciones, e identificar las influencias técnicas, operativas y climáticas más relevantes en la continuidad del servicio. Este desarrollo constituye una contribución significativa para fortalecer la toma de decisiones basada en evidencia, optimizar los planes de mantenimiento, apoyar el cumplimiento regulatorio y ofrecer un caso académico aplicado sobre el uso de técnicas analíticas en el sector eléctrico regulado, beneficiando a más de 600{.}000 usuarios.

El documento se organiza de la siguiente manera. En la Sección~\ref{sec:problema} se presenta el planteamiento y la formulación del problema de continuidad del servicio. En la Sección~\ref{sec:objetivos} se exponen el objetivo general y los objetivos específicos, alineados con las preguntas de investigación. En la Sección~\ref{sec:marco-referencia} se desarrolla el marco teórico (definiciones regulatorias de SAIDI y SAIFI, exclusiones y factores asociados) y los antecedentes sobre modelación de interrupciones e indicadores de confiabilidad. En la Sección~\ref{sec:datos} se describe la fuente y la estructura de los datos empleados. En la Sección~\ref{sec:metodologia} se detallan las etapas metodológicas ---desde la extracción e integración hasta el EDA, la ingeniería de variables y el modelamiento--- orientadas al pronóstico de la frecuencia y el tiempo de duración de interrupciones eléctricas. Posteriormente, en la Sección~\ref{sec:resultados} se presentan los resultados por etapa; en la Sección~\ref{sec:discusion} se discute el alcance, los aportes y las limitaciones frente a cada objetivo; y en la Sección~\ref{sec:conclusion} se consolidan las conclusiones del trabajo.
